\chapter{結論與未來工作}
\label{ch:conclusion-zh}

\section{結論}
\label{sec:conclude-main-zh}

本論文始於一個具體之經驗異常：於本文所用之稀疏評論衍生知識圖譜上，四個代表性 KG-aware 推薦器——KGAT、KGCL、MCCLK、KGRec——均一致地敗給純 LightGCN（第~\ref{ch:intro-zh}~章，並於第~\ref{sec:eval-main-zh}~節重現）。我們所主張之共通架構成因為：這些 baseline 將 KG 配置為計分管線中之強制元件，KG 嵌入參與訊息傳遞，且模型結構上不存在可在 KG 不可靠時將其脫接之開關。原本期望以語意訊號姿態被收割之噪音，因此進入產生協同表示之同一梯度流；當 KG 品質低於某門檻時，此交易即轉為淨負。

我們的回應 RA-GARK（第~\ref{ch:methodology-zh}~章）圍繞單一設計原則組織（第~\ref{sec:approach-principle-zh}~節）：KG 應作為一條\emph{旁路通道}，其對最終分數之貢獻由一可學習閘決定，此閘於初始化即偏向安全預設（純協同過濾），僅在訓練訊號指示 KG 通道能降低損失時才被打開。三項具體設計選擇落實此原則，並於第~\ref{sec:eval-ablation-zh}~節逐一衡量：

\begin{itemize}
  \item \textbf{KG-SVD 初始化}逐物品面向槽，於訓練起點即保留面向誘導語意子空間之角度幾何。移除其代價為 $5.3\%$ NDCG@$20$。
  \item \textbf{Softmax 合理化遮罩}對 $A$ 個潛在槽進行歸一化，編碼「使用者對某物品之偏好由少數互相競爭之語意面向所主導」之假設。將跨槽 Softmax 換為元素級 sigmoid 之代價為 $7.0\%$ NDCG@$20$——是本論文中單一元件影響最大者。
  \item \textbf{區域偏置融合閘}，其純量偏置 $b = +5$ 之初始化使 $\alpha^{(0)} \approx 0.993$，RA-GARK 訓練起點本質上即為純 LightGCN。還原為 $b = 0$ 之代價為 $5.3\%$ NDCG@$20$。
\end{itemize}

三項齊備下，RA-GARK 於本基準上取得 NDCG@$20 = 0.1238$，較最強 KG-aware baseline（KGRec）提升 $+13.1\%$，較純 LightGCN 提升 $+5.0\%$，且每 epoch 牆鐘成本基本相當（第~\ref{sec:eval-cost-zh}~節）。其中 $+5.0\%$ 在方法論上更具意義：於同一筆資料上，四個既發表之 KG-aware 方法將 KG 訊號變為淨負債，RA-GARK 之設計選擇則將其翻回淨資產。

因此本論文之一句話總結為：\emph{當 KG 不可靠時，架構最需要的不是更好的 KG 聚合器，而是一個能將 KG 結構性脫接之開關}。

\section{方法論啟示：注意力歸一化}
\label{sec:conclude-insight-zh}

本論文觀察到之單一元件最大影響並不在 RA-GARK 所貢獻之任何新架構元件，而在於推薦文獻中常被視為調參細節之選擇：合理化模組之注意力頭採用 Softmax（跨面向競爭）或 sigmoid（逐面向獨立）。於本基準上，僅將此單一模組之 Softmax 替換為 sigmoid 即損失 $7.0\%$ NDCG@$20$（第~\ref{sec:ablation-softmax-zh}~節），相對於一般超參數調參之 $1$--$3\%$ 擺動為大。

推薦文獻對此選擇無共識。DIN~\cite{zhou2018din} 對歷史物品採用 per-target sigmoid 注意力；NAIS~\cite{he2018nais}、AFM~\cite{xiao2017afm} 對特徵互動採用 Softmax 歸一化；KGRec~\cite{yang2023kgrec} 採邊層級 Bernoulli 決策，結構上接近 sigmoid。此等選擇皆與其目標任務之語意內部一致。本文之結果暗示：在合理化模組之角色為「於 $A$ 個競爭槽中選出最能在此使用者下代表此物品者」時，正確之語意假設為競爭，正確之歸一化因此為 Softmax；sigmoid baseline 之退化方式與槽權重漂移至均勻平均一致——其情境下合理化模組之選擇性即被默默禁用。

我們將此定位為假說而非方法論主張。證據僅來自單一資料集，於其他資料集、其他合理化結構之任務上其效應量級未知。此假說具體至足以由複製驗證證偽：\emph{當合理化模組之語意為「自少數互相競爭之潛在元件中選擇」時，跨元件 Softmax 為承載性預設，將其替換為獨立 sigmoid 頭會以遠超一般超參數雜訊之幅度降低排序品質}。若此於我們的單一設定之外仍成立，後續合理化結構之設計即應顯式處理此選擇。

\section{限制}
\label{sec:conclude-limits-zh}

本論文之貢獻受三項限制所界定，我們明確陳述以避免主張範圍模糊。

\paragraph{單一稀疏基準。} 本論文所有量化結果皆於單一資料集上報告：Amazon Books 評論之一子集，含 $905$ 位使用者、$1{,}399$ 件物品，以及 $2{,}098$ 個面向標籤上之 $3{,}370$ 條物品—面向 KG 邊（第~\ref{sec:eval-dataset-zh}~節）。選擇此基準是因其展現本論文所動機之經驗異常，但並未主張 RA-GARK 之設計選擇可不變地轉移至其他推薦資料集或其他 KG 架構。特別是，核心觀察——純 LightGCN 勝過所測四個 KG-aware baseline——於 KG 較稠密、較經策劃、或與推薦訊號較對齊之資料集上未必成立。

\paragraph{KG 建構為前作。} 物品—面向 KG 本身由~\cite{ho2024reviewkg}~之流程建構，結合視覺與文字編碼器（ResNet、Qwen-VL、BART、Mistral）以抽取評論衍生面向出現。我們直接採用此 KG，不對建構流程提出修改。本論文之貢獻完全位於建模面，視 KG 為固定輸入。

\paragraph{未驗證稠密 KG 情境。} RA-GARK 為稀疏 KG 情境而設計並評估，此情境下 KG 訊號不可靠。我們不對其於 KG 稠密且準確時之行為作任何主張。在該情境——KGAT、KGIN 之傳統主場——「預設 KG 關閉」之偏置初始化閘正是錯誤之歸納先驗，而我們所批評之深度融合 baseline 反而可能為合適選擇。第~\ref{sec:conclude-future-zh}~節討論閘應如何重新配置以處理對稱情境。

\section{未來工作}
\label{sec:conclude-future-zh}

由以上貢獻與限制，自然引出四個方向。

\paragraph{注意力歸一化發現之多資料集複製。} 第~\ref{sec:conclude-insight-zh}~節所表述之方法論假說具體至足以於其他合理化感知推薦基準、其他物品屬性架構、以及相鄰領域之合理化模組（順序推薦、對話式推薦）上測試。一次乾淨複製——固定資料集與模型其餘部分、僅切換合理化頭之歸一化——將使其自單一資料集觀察提升為一般設計指引，或銳化其成立之條件。

\paragraph{稠密 KG 驗證。} 如限制所提，RA-GARK 尚未於 KG 稠密且策劃良好之資料集上評估。自然之後續為將 RA-GARK 與 KG-aware baseline 並列執行於標準稠密 KG 基準（如 Amazon-Book / Last-FM 配 Freebase 衍生 KG）。預期結果為區域偏置融合閘優雅退化為與純 LightGCN 加小幅 KG 貢獻相當之行為而不崩壞；至於 RA-GARK 是否於該情境匹敵深度融合 baseline，為一個值得回答之經驗問題。

\paragraph{閘控原則之跨領域轉移。} 區域偏置融合閘背後之架構原則——於初始化偏向安全預設、僅在訓練訊號指示輔助通道能降低損失時打開——並不專屬於 KG-aware 推薦。同樣構造原則上可轉移至任何主訊號穩健、輔助訊號高變異之情境：含側資訊之冷啟動、單一模態雜訊時之多模態融合、異質預測器之集成式後期融合。我們所觀察之核心效應是否能在此等轉移下倖存，為開放問題。

\paragraph{反向稀疏情境：CF 稀疏配 KG 稠密。} 本論文所研究情境之對稱反面，為使用者—物品互動圖稀疏（如冷啟動、新使用者、新領域）而 KG 稠密可靠之情況。彼時閘所應編碼之安全預設與本文所用者相反：LightGCN 此時為不可靠之一側，KG 衍生之全域視角應為預設打開之通道。具體而言，融合偏置應自 $+5$ 翻轉為負值，KG-SVD 之角色可能由初始化轉為主要推論路徑，區域視角則可能需要不同之正則化才能維持有用。此為 KGAT、KGIN 等稠密 KG 工作之傳統設定；本文分析暗示\emph{於初始化偏向安全預設}此項設計動作可跨兩情境轉移，而偏置之具體方向則由資料集決定。對反向情境之完整架構處理需重新設計區域視角與閘方向性，留待後續研究。
